\chapter{A Concrete Naming Certificate for $\PeanoKernelUp$}
\label{ch:concrete-instances-peano-kernel-namecert}
\origin{human}

Following Peano and classical numerical analysis, the $\PeanoKernelUp$ packet isolates the finite-window data by which a quadrature-error functional is read without importing a host limiting theorem. It sits after the finite mesh carrier of \autoref{ch:concrete-instances-finite-partition-mesh-namecert}, the tagged Riemann-sum window surface of \autoref{ch:concrete-instances-riemann-sum-gauge-namecert}, the Darboux comparison row of \autoref{ch:concrete-instances-darboux-sum-namecert}, the regular-rational readback of \autoref{ch:concrete-instances-regseqrat-namecert}, the stream-window naming route of \autoref{ch:concrete-instances-streamname-namecert}, the dyadic tolerance ledger of \autoref{ch:concrete-instances-dyadicratcore-namecert}, and the terminal Real seal of \autoref{ch:concrete-instances-real-namecert}. The packet names only a displayed finite interval source, a mesh, a finite annihilation ledger for test polynomials, a kernel row, a sum consumer, regular readback, stream windows, dyadic tolerance, transport, provenance, and local naming; it does not claim host quadrature convergence, selected mesh limits, quotient equality of functions, measure completion, or ambient $\RealUp$ completeness.

\begin{definition}[Peano-kernel carrier]
\label{def:peano-kernel-carrier}
A $\PeanoKernelUp$ carrier is a finite $\BHist$ packet
$$
K=(I,M,A,Q,S,R,W,D,E,H,C,P,N)
$$
with the following displayed rows:
\begin{enumerate}
\item $I$ is the $\RealUp$ interval-source row naming the two endpoint seals, the oriented finite window, and the real-valued functional being inspected.
\item $M$ is the $\FinitePartitionMeshUp$ row, either a finite dyadic mesh or a finite uniform mesh, with displayed mesh points and adjacent-cell incidence.
\item $A$ is the test-polynomial annihilation ledger. It records the finite degree bound, the listed polynomial tests, and the zero-error verdict for those tests.
\item $Q$ is the kernel row. It stores the signed finite-window weights or cellwise kernel values used by the quadrature-error read after $I,M,A$ have been displayed.
\item $S$ is the Riemann/Darboux sum consumer row, read through $\RiemannSumGaugeUp$ or $\DarbouxSumUp$ data attached to the same mesh.
\item $R$ is the $\RegSeqRatUp$ readback row exposing the displayed finite error value or bound as a regular rational name.
\item $W$ is the $\StreamNameUp$ window row recording the finite read windows through which consumers may request the error value.
\item $D$ is the $\DyadicRatCoreUp$ tolerance ledger for the requested precision, comparison slack, and cellwise error budget.
\item $E$ is the terminal $\RealUp$ error seal reached only after $M,Q,S,R,W,D$ are public.
\item $H$ records componentwise $\hsame$ transport for interval source, mesh, annihilation ledger, kernel row, sum consumer, readback, stream windows, dyadic tolerance, Real seal, replay, provenance, and local naming.
\item $C$ records displayed $\Cont$ replay from every quadrature-error request through $I,M,A,Q,S,R,W,D,E,H,P,N$.
\item $P$ is the $\Pkg$ provenance over $\PeanoKernelUp$, $\RealUp$, $\FinitePartitionMeshUp$, $\RiemannSumGaugeUp$, $\DarbouxSumUp$, $\RegSeqRatUp$, $\StreamNameUp$, $\DyadicRatCoreUp$, $\BHist$, $\hsame$, $\Cont$, and $\NameCert$.
\item $N$ is the local $\NameCert_{\PeanoKernelUp}$ row.
\end{enumerate}
The classifier compares accepted carriers only through $I,M,A,Q,S,R,W,D,E,H,C,P,N$. It has no coordinate for an infinite refinement process, selected mesh limit, quotient equality of functions, measure completion, host convergence theorem, or ambient $\RealUp$ completeness.
\end{definition}

\begin{theorem}[Peano-kernel NameCert obligations]
\label{thm:peano-kernel-namecert-obligations}
For every accepted $\PeanoKernelUp$ carrier
$$
K=(I,M,A,Q,S,R,W,D,E,H,C,P,N),
$$
the local $\NameCert_{\PeanoKernelUp}$ surface consists exactly of carrier admission for $K$; Real interval-source exposure through $I$; finite dyadic or uniform mesh exposure through $M$; test-polynomial annihilation accountability through $A$; kernel-row compatibility through $Q$; Riemann/Darboux sum-consumer compatibility through $S$; regular-rational readback through $R$; stream-window read discipline through $W$; dyadic tolerance exactness through $D$; terminal Real error sealing through $E$; componentwise $\hsame$ transport through $H$; displayed $\Cont$ replay through $C$; $\Pkg$ provenance through $P$; local naming through $N$; and nonescape excluding host quadrature convergence, selected mesh limits, quotient function equality, measure completion, detached error representatives, and ambient $\RealUp$ completeness.
\end{theorem}
\begin{proof}
Carrier admission is projection from \autoref{def:peano-kernel-carrier}. The interval, mesh, annihilation ledger, kernel row, and sum consumer are the analytic-facing coordinates $I,M,A,Q,S$.

The readback coordinates are ordered. A quadrature-error request first fixes the displayed mesh and polynomial-test ledger, then reads the kernel and sum rows, then exposes the regular rational row $R$ through the stream window row $W$ under the dyadic tolerance ledger $D$.

The terminal seal is $E$. It is downstream of the finite mesh, kernel, sum, readback, stream-window, and tolerance rows, so it exports a Real error name rather than a host convergence theorem.

Finally $H,C,P,N$ transport, replay, source, and name exactly the same route. Since no carrier coordinate stores an infinite refinement process, selected mesh limit, quotient equality of functions, measure completion, detached error representative, or ambient completeness, the listed obligations exhaust the local NameCert surface.
\end{proof}

\begin{theorem}[Peano-kernel error-functional route]
\label{thm:peano-kernel-error-functional-route}
Every quadrature-error read of an accepted $\PeanoKernelUp$ carrier factors through the displayed finite route
$$
\RealUp \longmapsto \FinitePartitionMeshUp \longmapsto \PeanoKernelUp \longmapsto (\RiemannSumGaugeUp \text{ or } \DarbouxSumUp) \longmapsto \RegSeqRatUp \longmapsto \StreamNameUp \longmapsto \DyadicRatCoreUp \longmapsto \RealUp .
$$
Equivalently, the read can inspect only the finite mesh row $M$, the kernel row $Q$, the sum row $S$, the regular readback row $R$, the stream-window row $W$, the dyadic tolerance ledger $D$, the Real error seal $E$, and the structural rows $H,C,P,N$ after the interval source and annihilation ledger have been displayed.
\end{theorem}
\begin{proof}
Start with an accepted carrier $K=(I,M,A,Q,S,R,W,D,E,H,C,P,N)$ and a quadrature-error read. The only interval source is $I$, and the only mesh source is the displayed finite row $M$.

The annihilation ledger $A$ binds the polynomial tests before the kernel is read. The row $Q$ then provides the cellwise kernel or weight data, and $S$ consumes it through the displayed Riemann or Darboux sum surface attached to the same mesh.

The read cannot leave the finite route after $S$. It reaches the regular-rational row $R$, the stream-window row $W$, and the dyadic tolerance ledger $D$ before the terminal Real error seal $E$ is exposed.

The rows $H,C,P,N$ preserve that order under transport, replay, provenance, and local naming. Because the carrier contains no infinite refinement process, selected mesh limit, quotient function equality, measure completion, detached error representative, or ambient completeness coordinate, every accepted quadrature-error read factors through the stated finite route.
\end{proof}

\closureat{PeanoKernelUp}{seedStr}

\begin{closurestatus}{\PeanoKernelUp}
  \constructivestory{A $\PeanoKernelUp$ packet is generated from a RealUp interval source, a finite dyadic or uniform partition mesh, a test-polynomial annihilation ledger, a kernel row, a Riemann/Darboux sum consumer, RegSeqRatUp readback, StreamNameUp windows, a DyadicRatCoreUp tolerance ledger, a RealUp error seal, componentwise hsame transport, displayed Cont replay, Pkg provenance, and a local NameCert row.}
  \theoryclosure{\seedClosure}
  \scopeclosed{\autoref{def:peano-kernel-carrier}, \autoref{thm:peano-kernel-namecert-obligations}, and \autoref{thm:peano-kernel-error-functional-route} seed the finite quadrature-error route from interval source and finite mesh to kernel, sum, regular readback, stream windows, dyadic tolerance, and the Real error seal.}
  \formalstatus{\unformalizedV}
  \bridgestatus{none}
  \notclaimed{This chapter does not close host quadrature convergence, selected mesh limits, quotient function equality, measure completion, ambient $\RealUp$ completeness, Lean theorem checking, or bridge checking.}
  \upgradepath{Obligation closure requires checked carrier admission, interval-source exposure, finite mesh exposure, annihilation-ledger accountability, kernel-row compatibility, sum-consumer compatibility, regular readback, stream-window discipline, dyadic tolerance exactness, Real error sealing, transport, replay, provenance, and local naming for BEDC.Derived.PeanoKernelUp.}
\end{closurestatus}
